\documentclass[../../main.tex]{subfiles}

\begin{document}

\section{带号测度: Hahn 与 Jordan 分解}\label{section:ch17-signed-measure}

观察到若 $\mu_{1}$ 和 $\mu_{2}$ 是定义在同一个可测空间 $(X,\mathcal{M})$ 上的测度, 则对正数 $\alpha$ 和 $\beta$, 通过对 $\mathcal{M}$ 中所有 $E$ 设
\begin{align*}
\mu_{3}(E)=\alpha\cdot\mu_{1}(E)+\beta\cdot\mu_{2}(E)
\end{align*}
我们可以定义 $X$ 上的新的测度 $\mu_{3}$. 事实证明, 考虑测度的线性组合但系数可以是负的是重要的. 若我们试图定义 $\mathcal{M}$ 上的集函数 $\nu$ 为对 $\mathcal{M}$ 中所有 $E$,
\begin{align*}
\nu(E)=\mu_{1}(E)-\mu_{2}(E)
\end{align*}
会发生什么呢?

第一件可能发生的事情是 $\nu$ 不总是非负. 此外, $\nu(E)$ 甚至对使得 $\mu_{1}(E)=\mu_{2}(E)=\infty$ 的 $E\in\mathcal{M}$ 没有定义. 有了这些考虑我们给出以下定义.

\begin{definition}[带号测度]\label{definition:ch17-signed-measure}
谈到可测空间 $(X,\mathcal{M})$ 上的带号测度 $\nu$, 我们指的是具有以下性质的扩充实值集函数 $\nu:\mathcal{M}\to[-\infty,\infty]$:
\begin{description}
\item[（i）] $\nu$ 至多取到 $+\infty$, $-\infty$ 中的一个值.
\item[（ii）] $\nu(\varnothing)=0$.
\item[（iii）] 对任何不交可测集的可数族 $\{E_{k}\}_{k=1}^{\infty}$,
\begin{align*}
\nu\left(\bigcup_{k=1}^{\infty}E_{k}\right)=\sum_{k=1}^{\infty}\nu(E_{k})
\end{align*}
其中若 $\nu\left(\bigcup_{k=1}^{\infty}E_{k}\right)$ 是有限的, 级数 $\sum_{k=1}^{\infty}\nu(E_{k})$ 绝对收敛.
\end{description}
\end{definition}

测度是带号测度的特殊情形. 不难看出两个测度的差是一个带号测度. 事实上, 接下来的 Jordan 分解定理将告诉我们, 每个带号测度就是两个这样的测度的差.

令 $\nu$ 为带号测度. 我们说一个集合 $A$ 是正的(关于 $\nu$), 若 $A$ 是可测的且对 $A$ 的每个可测子集 $E$ 有 $\nu(E)\geqslant 0$. $\nu$ 在正集的可测子集上的限制是一个测度. 类似地, 称一个集合 $B$ 是负的(关于 $\nu$), 若 $B$ 是可测的且 $B$ 的每个可测子集 $E$ 有非正的 $\nu$ 测度. $-\nu$ 在负集的可测子集上的限制也是一个测度. 可测集称为是关于 $\nu$ 的零集, 若它的每个可测子集的 $\nu$ 测度为零. 读者必须仔细区分零集与测度为零的集合: 每个零集的测度必须为零, 测度为零的集合可以是两个集合的并, 它们的测度不为零但互为相反数. 根据测度的单调性, 一个集合关于某个测度是零当且仅当它的测度为零. 由于一个带号测度 $\nu$ 不取 $+\infty$ 和 $-\infty$ 值, 对可测集 $A$ 和 $B$,
\begin{align}
\text{若 }A\subseteq B\text{ 且 }|\nu(B)|<\infty,\quad\text{则 }|\nu(A)|<\infty.\label{align:ch17-finite-monotonicity}
\end{align}

\begin{proposition}\label{proposition:ch17-positive-set}
令 $\nu$ 为可测空间 $(X,\mathcal{M})$ 上的带号测度. 则正集的每个可测子集的测度是正的且可数个正集的并的测度是正的.
\end{proposition}

\begin{proof}
由正集的定义可知第一个陈述平凡成立. 为证明第二个陈述, 令 $A$ 为正集的可数族 $\{A_{k}\}_{k=1}^{\infty}$ 的并. 令 $E$ 为 $A$ 的可测子集. 定义 $E_{1}=E\cap A_{1}$. 对于 $k\geqslant 2$, 定义
\begin{align*}
E_{k}=[E\cap A_{k}]\sim[A_{1}\cup\cdots\cup A_{k-1}].
\end{align*}
则每个 $E_{k}$ 是正集 $A_{k}$ 的可测子集, 因此 $\nu(E_{k})\geqslant 0$. 由于 $E$ 是可数不交族 $\{E_{k}\}_{k=1}^{\infty}$ 的并,
\begin{align*}
\nu(E)=\sum_{k=1}^{\infty}\nu(E_{k})\geqslant 0.
\end{align*}
因此 $A$ 是一个正集.
\end{proof}

\begin{lemma}[Hahn 引理]\label{lemma:ch17-hahn}
令 $\nu$ 为可测空间 $(X,\mathcal{M})$ 上的带号测度, 而 $E$ 是满足 $0<\nu(E)<\infty$ 的可测集. 则存在 $E$ 的正的且有正测度的可测子集 $A$.
\end{lemma}

\begin{proof}
若 $E$ 本身是正集, 则引理得证. 否则, $E$ 包含负测度集. 令 $m_{1}$ 为使得存在测度小于 $-1/m_{1}$ 的可测集的最小自然数. 选取可测集 $E_{1}\subseteq E$ 使得 $\nu(E_{1})<-1/m_{1}$. 令 $n$ 为自然数, 其中满足以下条件的自然数 $m_{1},\cdots,m_{n}$ 与可测集 $E_{1},\cdots,E_{n}$ 已被选取: 对于 $1\leqslant k\leqslant n$, $m_{k}$ 是使得 $E\sim\bigcup_{j=1}^{k-1}E_{j}$ 存在测度小于 $-1/m_{k}$ 的可测子集的最小自然数, 而 $E_{k}$ 是 $\left[E\sim\bigcup_{j=1}^{k-1}E_{j}\right]$ 的满足 $\nu(E_{k})<-1/m_{k}$ 的子集.

若该选取过程结束, 则完成证明. 否则, 定义
\begin{align*}
A=E\sim\bigcup_{k=1}^{\infty}E_{k},\quad\text{使得 }E=A\cup\left[\bigcup_{k=1}^{\infty}E_{k}\right]\text{ 是 }E\text{ 的一个不交分解}.
\end{align*}
由于 $\bigcup_{k=1}^{\infty}E_{k}$ 是 $E$ 的可测子集且 $|\nu(E)|<\infty$, 由\eqref{align:ch17-finite-monotonicity}和 $\nu$ 的可数可加性,
\begin{align*}
-\infty<\nu\left(\bigcup_{k=1}^{\infty}E_{k}\right)=\sum_{k=1}^{\infty}\nu(E_{k})\leqslant\sum_{k=1}^{\infty}-1/m_{k}.
\end{align*}
因此 $\lim_{k\to\infty}m_{k}=\infty$. 我们宣称 $A$ 是一个正集. 事实上, 若 $B$ 是 $A$ 的可测子集, 则对每个 $k$,
\begin{align*}
B\subseteq A\subseteq E\sim\left[\bigcup_{j=1}^{k-1}E_{j}\right].
\end{align*}
由 $m_{k}$ 的最小选取, $\nu(B)\geqslant-1/(m_{k}-1)$. 由于 $\lim_{k\to\infty}m_{k}=\infty$, 有 $\nu(B)\geqslant 0$. 因此 $A$ 是一个正集. 剩下来仅须证明 $\nu(A)>0$. 这可从 $\nu$ 的有限可加性得出, 因为 $\nu(E)>0$ 且 $\nu(E\sim A)=\nu\left(\bigcup_{k=1}^{\infty}E_{k}\right)=\sum_{k=1}^{\infty}\nu(E_{k})<0$.
\end{proof}

\begin{theorem}[Hahn 分解定理]\label{theorem:ch17-hahn-decomposition}
令 $\nu$ 为可测空间 $(X,\mathcal{M})$ 上的带号测度. 则存在关于 $\nu$ 的正集 $A$ 和关于 $\nu$ 的负集 $B$, 使得
\begin{align*}
X=A\cup B\text{ 且 }A\cap B=\varnothing.
\end{align*}
\end{theorem}

\begin{proof}
不失一般性, 我们假设 $+\infty$ 是 $\nu$ 取不到的无穷值. 令 $\mathcal{P}$ 为 $X$ 的正子集族且定义
\begin{align*}
\lambda=\sup\{\nu(E)\mid E\in\mathcal{P}\}.
\end{align*}
则 $\lambda\geqslant 0$, 因为 $\mathcal{P}$ 包含空集. 令 $\{A_{k}\}_{k=1}^{\infty}$ 为正集的可数族使得 $\lambda=\lim_{k\to\infty}\nu(A_{k})$. 定义 $A=\bigcup_{k=1}^{\infty}A_{k}$, 根据命题\cref{proposition:ch17-positive-set}, 集合 $A$ 是一个正集, 因而 $\lambda\geqslant\nu(A)$. 另一方面, 对每个 $k$, $A\sim A_{k}\subseteq A$, 因而 $\nu(A\sim A_{k})\geqslant 0$. 因此,
\begin{align*}
\nu(A)=\nu(A_{k})+\nu(A\sim A_{k})\geqslant\nu(A_{k}).
\end{align*}
于是 $\nu(A)\geqslant\lambda$. 因此 $\nu(A)=\lambda$ 且 $\lambda<\infty$ 由于 $\nu$ 取不到 $\infty$ 值.

令 $B=X\sim A$. 我们用反证法证明 $B$ 是负的. 假设 $B$ 不是负的. 则存在 $B$ 的具有正测度的子集 $E$, 根据 Hahn 引理, $B$ 的子集 $E_{0}$ 既是正的又有正的测度. 则 $A\cup E_{0}$ 是正集且
\begin{align*}
\nu(A\cup E_{0})=\nu(A)+\nu(E_{0})>\lambda,
\end{align*}
这与 $\lambda$ 的选取矛盾.
\end{proof}

将 $X$ 分解为两个不交集合 $A$ 和 $B$ 的并, 其中 $A$ 关于 $\nu$ 是正的, 而 $B$ 是负的, 称为关于 $\nu$ 的 Hahn 分解. 前一个定理告诉我们对每个带号测度的 Hahn 分解的存在性. 这样的分解可以不唯一. 的确, 若 $\{A,B\}$ 是关于 $\nu$ 的 Hahn 分解, 则通过从 $A$ 中去掉一个零集 $E$ 且将它移到 $B$ 上, 我们得到另一个 Hahn 分解 $\{A\sim E,B\cup E\}$.

若 $\{A,B\}$ 是关于 $\nu$ 的 Hahn 分解, 通过设
\begin{align*}
\nu^{+}(E)=\nu(E\cap A)\text{ 与 }\nu^{-}(E)=\nu(E\cap B)
\end{align*}
我们定义了两个测度 $\nu^{+}$ 和 $\nu^{-}$. 它们满足 $\nu=\nu^{+}-\nu^{-}$. $(X,\mathcal{M})$ 上的两个测度 $\nu_{1}$ 和 $\nu_{2}$ 称为相互奇异的(记为 $\nu_{1}\perp\nu_{2}$), 若存在不交可测集 $A$ 和 $B$ 使得 $X=A\cup B$, $\nu_{1}(A)=\nu_{2}(B)=0$. 上面定义的测度 $\nu^{+}$ 和 $\nu^{-}$ 是相互奇异的. 因此我们已证明了以下命题的存在性部分. 唯一性部分的证明留给读者(见习题 13).

\begin{theorem}[Jordan 分解定理]\label{theorem:ch17-jordan-decomposition}
令 $\nu$ 为可测空间 $(X,\mathcal{M})$ 上的带号测度. 则 $(X,\mathcal{M})$ 上存在两个相互奇异的测度 $\nu^{+}$ 和 $\nu^{-}$ 满足 $\nu=\nu^{+}-\nu^{-}$. 此外, 仅存在一对这样的相互奇异的测度.
\end{theorem}

该定理给出的带号测度 $\nu$ 的分解称为 $\nu$ 的 Jordan 分解. 测度 $\nu^{+}$ 和 $\nu^{-}$ 分别称为 $\nu$ 的正部(变差)与负部(变差). 由于 $\nu$ 至多取到值 $+\infty$ 和 $-\infty$ 中的一个, $\nu^{+}$ 或 $\nu^{-}$ 必须是有限的. 若它们都是有限的, 我们称 $\nu$ 为有限带号测度. $\mathcal{M}$ 上的测度 $|\nu|$ 定义为对所有 $E\in\mathcal{M}$, $|\nu|(E)=\nu^{+}(E)+\nu^{-}(E)$.

我们把以下等式的证明留给读者作为练习.
\begin{align}
|\nu|(X)=\sup\sum_{k=1}^{n}|\nu(E_{k})|\label{align:ch17-total-variation}
\end{align}
其中上确界对 $X$ 的可测子集的所有有限不交族 $\{E_{k}\}_{k=1}^{n}$ 取, 因此 $|\nu|(X)$ 称为 $\nu$ 的全变差, 记为 $\|\nu\|_{\mathrm{var}}$.

\begin{example}\label{example:ch17-signed-measure-integral}
令 $f:\mathbf{R}\to\mathbf{R}$ 为 $\mathbf{R}$ 上的 Lebesgue 可积函数. 对于 Lebesgue 可测集 $E$, 定义 $\nu(E)=\int_{E}f\,dm$. 我们从积分的可数可加性(4.5 节定理 20)推出, $\nu$ 是可测空间 $(\mathbf{R},\mathcal{L})$ 上的带号测度. 定义 $A=\{x\in\mathbf{R}\mid f(x)\geqslant 0\}$ 和 $B=\{x\in\mathbf{R}\mid f(x)<0\}$, 且对 Lebesgue 可测集 $E$, 定义
\begin{align*}
\nu^{+}(E)=\int_{A\cap E}f\,dm\text{ 与 }\nu^{-}(E)=-\int_{B\cap E}f\,dm.
\end{align*}
则 $\{A,B\}$ 是 $\mathbf{R}$ 关于带号测度 $\nu$ 的 Hahn 分解. 此外, $\nu=\nu^{+}-\nu^{-}$ 是 $\nu$ 的 Jordan 分解.
\end{example}

\begin{exercise}\label{exercise:ch17-12}
在上面的例子中, 令 $E$ 为满足 $0<\nu(E)<\infty$ 的 Lebesgue 可测集. 找到包含于 $E$ 的正集 $A$ 使得 $\nu(A)>0$.
\end{exercise}

\begin{exercise}\label{exercise:ch17-13}
令 $\mu$ 为测度而 $\mu_{1}$ 和 $\mu_{2}$ 是测度空间 $(X,\mu)$ 上使得 $\mu=\mu_{1}-\mu_{2}$ 的相互奇异的测度. 证明 $\mu_{2}=0$. 用此证明 Jordan 分解定理的唯一性断言.
\end{exercise}

\begin{exercise}\label{exercise:ch17-14}
证明: 若 $E$ 是任何可测集, 则
\begin{align*}
-\nu^{-}(E)\leqslant\nu(E)\leqslant\nu^{+}(E)\text{ 且 }|\nu(E)|\leqslant|\nu|(E).
\end{align*}
\end{exercise}

\begin{exercise}\label{exercise:ch17-15}
证明: 若 $\nu_{1}$ 和 $\nu_{2}$ 是任何两个有限带号测度, 则 $\alpha\nu_{1}+\beta\nu_{2}$ 也是如此, 其中 $\alpha$ 和 $\beta$ 是实数. 证明
\begin{align*}
|\alpha\nu|=|\alpha||\nu|\text{ 且 }|\nu_{1}+\nu_{2}|\leqslant|\nu_{1}|+|\nu_{2}|,
\end{align*}
其中 $\nu\leqslant\mu$ 意味着对所有可测集 $E$, $\nu(E)\leqslant\mu(E)$.
\end{exercise}

\begin{exercise}\label{exercise:ch17-16}
证明\eqref{align:ch17-total-variation}.
\end{exercise}

\begin{exercise}\label{exercise:ch17-17}
令 $\mu$ 和 $\nu$ 为有限带号测度. 定义 $\mu\wedge\nu=\frac{1}{2}(\mu+\nu-|\mu-\nu|)$ 与 $\mu\vee\nu=\mu+\nu-\mu\wedge\nu$.
\begin{description}
\item[（i）] 证明带号测度 $\mu\wedge\nu$ 比 $\mu$ 和 $\nu$ 小但比任何小于 $\mu$ 和 $\nu$ 的其他带号测度大.
\item[（ii）] 证明带号测度 $\mu\vee\nu$ 大于 $\mu$ 和 $\nu$ 但小于任何大于 $\mu$ 和 $\nu$ 的其他测度.
\item[（iii）] 若 $\mu$ 和 $\nu$ 是正测度, 证明它们相互奇异当且仅当 $\mu\wedge\nu=0$.
\end{description}
\end{exercise}

\end{document}